
\section{Problem 616 (\#616)}

\subsection*{1) ROUND-2 OBJECTIVE}

Path (C): identify a minimal formal correction to the Round-1 restatement, and then use it to obtain exact information for small uniformities. Concretely, I will replace the Round-1 ``induced subhypergraph'' phrasing by the original subgraph/edge-subfamily formulation, and then prove the exact values
\[
t(3)=t(4)=t(5)=1,\qquad t(6)=2,
\]
together with a general improvement of the Round-1 bound from $\tau(G)\le r$ to $\tau(G)\le r-1$.

\subsection*{2) Round-1 FOUNDATION USED}

I use the following Round-1 consequences of the local hypothesis.

\begin{enumerate}
\item For $r\ge 3$, any two edges of $G$ intersect.
\item For $r\ge 3$, any three edges of $G$ intersect.
\item For $r\ge 3$, for any two edges $e,f$, the set $e\cap f$ is a transversal of $G$.
\end{enumerate}

I will not re-prove these; I use them as vetted Round-1 input.

\subsection*{3) NEW INSIGHT / TOOL (ROUND-2)}

The new ingredient is a ``minimal non-Helly subfamily'' lemma.

If $\mathcal{F}=\{E_1,\dots,E_m\}$ is an $r$-uniform family with empty total intersection and minimal with respect to that property, then
\[
m\le r+1
\quad\text{and}\quad
\left|\bigcup_{i=1}^m E_i\right|\le m(r-m+2).
\]

This turns the local threshold $3r-3$ into a concrete obstruction for small $r$.

A second new observation is that if $T$ is a minimal transversal of size at least $2$, then two witness edges for distinct vertices of $T$ intersect in at most $r-1$ vertices; since every pair-intersection is a transversal by Round~1, this yields the global improvement
\[
\tau(G)\le r-1.
\]

\subsection*{4) ATTACK PLAN (ROUND-2)}

The exact gap left by Round~1 was this: the local condition forced $3$-wise intersection, but there was no mechanism turning that into either (i) exact small-$r$ values or (ii) a sharper global upper bound than $\tau(G)\le r$.

The plan is:

\begin{enumerate}
\item reformulate the original hypothesis exactly in terms of edge-subfamilies;
\item study a minimal edge-subfamily with empty total intersection;
\item bound the union size of such a minimal obstruction;
\item compare that bound with $3r-3$.
\end{enumerate}

For $r\le 5$, this forces a contradiction, giving $\tau(G)=1$. For $r=6$, only a $4$-edge obstruction survives; that obstruction has enough rigid structure to force a $2$-vertex transversal.

\subsection*{5) WORK (ROUND-2)}

\paragraph{Corrected local formulation.}
The original problem is about \emph{all subgraphs} on at most $3r-3$ vertices, not just induced subgraphs. The exact equivalent formulation is:

\medskip
\noindent
\emph{For every nonempty edge-subfamily $\mathcal{H}\subseteq E(G)$ with}
\[
\left|\bigcup \mathcal{H}\right|\le 3r-3,
\]
\emph{we have}
\[
\bigcap_{e\in \mathcal{H}} e\neq \varnothing.
\]

Indeed, such an edge-subfamily spans a subgraph on $\left|\bigcup \mathcal{H}\right|$ vertices, and conversely every subgraph is determined by its edge set.

\paragraph{Lemma 5.1 (minimal non-Helly subfamily).}
Let $\mathcal{F}=\{E_1,\dots,E_m\}$ be an $r$-uniform family such that
\[
\bigcap_{i=1}^m E_i=\varnothing,
\]
but every proper subfamily has nonempty intersection. Then:

\begin{enumerate}
\item $m\le r+1$;
\item $\left|\bigcup_{i=1}^m E_i\right|\le m(r-m+2)$.
\end{enumerate}

\paragraph{Proof.}
For each $i$, minimality gives a vertex
\[
x_i\in \bigcap_{j\ne i} E_j\setminus E_i.
\]
The vertices $x_1,\dots,x_m$ are distinct: if $x_i=x_j$ with $i\ne j$, then that common vertex would belong to $E_j$ (because $x_i\in \cap_{k\ne i} E_k$) and would also fail to belong to $E_j$ (because $x_j\notin E_j$), a contradiction.

Now fix $i$. Since $x_j\in E_i$ for every $j\ne i$, the edge $E_i$ contains at least the $m-1$ distinct vertices
\[
\{x_j:\ j\ne i\}.
\]
As $|E_i|=r$, this implies $m-1\le r$, so $m\le r+1$.

Let
\[
X=\{x_1,\dots,x_m\}.
\]
Each $E_i$ already contains the $m-1$ vertices $X\setminus\{x_i\}$, so it contributes at most
\[
r-(m-1)=r-m+1
\]
additional vertices outside $X$. Therefore
\[
\left|\bigcup_{i=1}^m E_i\right|
\le |X|+\sum_{i=1}^m \bigl(r-m+1\bigr)
= m+m(r-m+1)
= m(r-m+2).
\]
This proves the lemma.
\hfill $\square$

\paragraph{Corollary 5.2.}
If $r\in\{3,4,5\}$ and $G$ satisfies the local hypothesis, then all edges of $G$ have a common vertex. Equivalently,
\[
t(3)=t(4)=t(5)=1.
\]

\paragraph{Proof.}
Assume for contradiction that $G$ has no common vertex. By Round~1, every three edges intersect. Hence any minimal edge-subfamily with empty total intersection has size
\[
m\ge 4.
\]
By Lemma~5.1, necessarily $m\le r+1$ and its union has size at most $m(r-m+2)$.

Now check the finitely many cases.

\begin{itemize}
\item If $r=3$, then necessarily $m=4$, and
\[
m(r-m+2)=4(3-4+2)=4\le 6=3r-3.
\]
\item If $r=4$, then $m\in\{4,5\}$, and
\[
4(4-4+2)=8\le 9,\qquad 5(4-5+2)=5\le 9.
\]
\item If $r=5$, then $m\in\{4,5,6\}$, and
\[
4(5-4+2)=12=3r-3,\qquad 5(5-5+2)=10\le 12,\qquad 6(5-6+2)=6\le 12.
\]
\end{itemize}

So in every case the minimal obstruction has union size at most $3r-3$. By the local hypothesis, such a subfamily must have nonempty intersection, contradicting the choice of the obstruction.

Therefore all edges of $G$ intersect, so $\tau(G)=1$.
\hfill $\square$

\paragraph{Proposition 5.3 (general improvement of the Round-1 upper bound).}
For every $r\ge 3$, if $G$ satisfies the local hypothesis, then
\[
\tau(G)\le r-1.
\]

\paragraph{Proof.}
If $\tau(G)=1$ there is nothing to prove. Assume $\tau(G)\ge 2$, and let $T$ be a minimal transversal of size $\tau(G)$.

For each $x\in T$, choose a witness edge $E_x$ with
\[
E_x\cap T=\{x\}.
\]
Pick two distinct vertices $x,y\in T$. By the Round-1 foundation, the set
\[
E_x\cap E_y
\]
is a transversal of $G$, so
\[
\tau(G)\le |E_x\cap E_y|.
\]
On the other hand, because $E_x\cap T=\{x\}$ and $E_y\cap T=\{y\}$ with $x\ne y$, the intersection $E_x\cap E_y$ contains no vertex of $T$, hence
\[
|E_x\cap E_y|\le r-1.
\]
Therefore $\tau(G)\le r-1$.
\hfill $\square$

\paragraph{Construction 5.4 (a $2$-cover example for every $r\ge 6$).}
Let
\[
X=\{x_1,x_2,x_3,x_4\},
\]
and let $A_1,A_2,A_3,A_4$ be pairwise disjoint sets of size $r-3$, all disjoint from $X$.
Define four $r$-edges by
\[
E_i=(X\setminus \{x_i\})\cup A_i
\qquad (1\le i\le 4).
\]

Then:

\begin{enumerate}
\item every three of the four edges intersect;
\item the total intersection is empty;
\item the local hypothesis holds whenever $r\ge 6$;
\item $\tau(\{E_1,E_2,E_3,E_4\})=2$.
\end{enumerate}

\paragraph{Verification.}
Any three edges contain the unique witness vertex omitted by the fourth. The total intersection is empty because no $x_i$ belongs to $E_i$, and the $A_i$ are disjoint.

The union of all four edges has size
\[
|X|+\sum_{i=1}^4 |A_i|=4+4(r-3)=4r-8.
\]
For $r\ge 6$,
\[
4r-8>3r-3.
\]
So a subfamily on at most $3r-3$ vertices can contain at most three of the four edges; every such subfamily has nonempty intersection. Hence the local condition holds.

Finally, no single vertex hits all four edges, so $\tau\ge 2$. But if $a\in A_1$, then the $2$-set $\{x_1,a\}$ hits $E_1$ via $a$ and hits $E_2,E_3,E_4$ via $x_1$. Hence $\tau=2$.

Thus
\[
t(r)\ge 2 \qquad\text{for all } r\ge 6.
\]
\hfill $\square$

\paragraph{Theorem 5.5.}
\[
t(6)=2.
\]

\paragraph{Proof.}
Construction~5.4 gives $t(6)\ge 2$, so it remains to prove $t(6)\le 2$.

Let $G$ be a $6$-uniform hypergraph satisfying the local hypothesis. If all edges have a common vertex, then $\tau(G)=1$. So assume $G$ has no common vertex.

By Round~1, every three edges intersect. Choose a minimal edge-subfamily
\[
\mathcal{F}=\{E_1,\dots,E_m\}
\]
with empty total intersection. Then $m\ge 4$.

By Lemma~5.1,
\[
\left|\bigcup \mathcal{F}\right|\le m(8-m).
\]
If $m=5,6,$ or $7$, then
\[
5(8-5)=15,\qquad 6(8-6)=12,\qquad 7(8-7)=7,
\]
all of which are at most $15=3r-3$. That would contradict the local hypothesis. Hence the only possibility is
\[
m=4.
\]

Now apply Lemma~5.1 to this minimal $4$-edge obstruction. Choose witnesses
\[
x_i\in \bigcap_{j\ne i} E_j\setminus E_i
\qquad (i=1,2,3,4),
\]
and let
\[
X=\{x_1,x_2,x_3,x_4\}.
\]
Each $E_i$ contains $X\setminus\{x_i\}$ and has size $6$, so it has exactly three additional vertices outside $X$. Therefore
\[
\left|\bigcup_{i=1}^4 E_i\right|\le 4+4\cdot 3=16.
\]
Because $\mathcal{F}$ has empty intersection, the local hypothesis forces
\[
\left|\bigcup_{i=1}^4 E_i\right|>15.
\]
Hence the union has size exactly $16$, so the upper bound is tight. Therefore the three extra vertices of each $E_i$ are disjoint from $X$ and pairwise disjoint across $i=1,2,3,4$. In particular,
\[
E_i\cap E_j = X\setminus\{x_i,x_j\}
\qquad (i\ne j).
\]

Now let $F$ be any edge of $G$. By the Round-1 foundation, any three edges intersect, so for every pair $i\ne j$,
\[
F\cap E_i\cap E_j\ne \varnothing.
\]
Since $E_i\cap E_j=X\setminus\{x_i,x_j\}$, it follows that
\[
F\cap (X\setminus\{x_i,x_j\})\ne \varnothing
\qquad\text{for every } i\ne j.
\]
Equivalently, the set $F\cap X$ meets every $2$-subset of $X$. This is possible only if
\[
|F\cap X|\ge 3.
\]
Indeed, if $|F\cap X|\le 2$, then the complement $X\setminus (F\cap X)$ contains a $2$-subset disjoint from $F\cap X$.

Fix the pair
\[
P=\{x_1,x_2\}.
\]
Each of the four core edges $E_1,E_2,E_3,E_4$ meets $P$, because each $E_i$ contains three of the four vertices of $X$. Any other edge $F$ also meets $P$, because $F\cap X$ has size at least $3$, and every $3$-subset of a $4$-set meets every fixed $2$-subset.

Thus $P$ is a transversal of $G$, so $\tau(G)\le 2$.

Combining this with the lower construction gives
\[
t(6)=2.
\]
\hfill $\square$

\subsection*{6) ADVERSARIAL VERIFICATION}

\begin{enumerate}
\item The only formal correction made to Round~1 is replacing the stronger ``induced subhypergraph'' wording by the original subgraph/edge-subfamily formulation. All Round-1 lemmas used here remain valid under the original statement.
\item In Lemma~5.1, the witness vertices $x_i$ are indeed distinct; this is the only subtle point in the proof.
\item In Corollary~5.2, the contradiction uses only the existence of \emph{one} minimal empty-intersection subfamily with union at most $3r-3$, so the local hypothesis applies exactly as required.
\item In Construction~5.4, every potentially problematic subfamily is either:
\begin{itemize}
\item a subfamily of at most three edges, which has nonempty intersection, or
\item the full $4$-edge family, whose union has size $4r-8>3r-3$.
\end{itemize}
So the local condition is verified.
\item In the proof of $t(6)=2$, the step forcing $|F\cap X|\ge 3$ uses only the exact form of the six pair-intersections
\[
E_i\cap E_j=X\setminus\{x_i,x_j\},
\]
which was justified by the equality case in the union bound.
\end{enumerate}

\subsection*{7) FINAL}

\textbf{UNRESOLVED (BUT STRICTLY ADVANCED).}

The main new results of this round are:

\begin{enumerate}
\item the exact values
\[
t(3)=t(4)=t(5)=1,\qquad t(6)=2;
\]
\item the general lower construction
\[
t(r)\ge 2 \qquad (r\ge 6);
\]
\item the general upper bound
\[
\tau(G)\le r-1
\]
for every admissible $r$-uniform hypergraph.
\end{enumerate}

So the global problem is still open for $r\ge 7$, but the small cases are now settled through $r=6$, and the Round-1 bound $\tau(G)\le r$ has been sharpened to $\tau(G)\le r-1$.

\subsection*{8) COMPLETION ESTIMATE}

COMPLETION: 72\%

\subsection*{9) REFERENCES}

\begin{itemize}
\item T. F. Bloom, \emph{Erd\H{o}s Problem \#616}, accessed 2026-03-08.
\item Round-1 notes from the previous iteration of this investigation (used only for the three structural lemmas listed in Section~2).
\end{itemize}

